%=====================================================================
% REPRESENTACAO GRAFICA REVISADA - TOMADAS
%=====================================================================
\section{Leitura gráfica de tomadas}

\begin{frame}{Tomada 2P+T — representação por bornes}
\small
\begin{columns}[T,onlytextwidth]
\column{.52\textwidth}
\centering
\begin{tikzpicture}[scale=.88,every node/.style={font=\scriptsize}]
\draw[corNorma,fill=corNorma!5,rounded corners](-2.2,-1.6)rectangle(2.2,1.6);
\node at(0,1.18){T1 — 2P+T};
\node[circle,draw=corPrim,minimum size=5pt,inner sep=0pt](a)at(-1.05,.2){};
\node[circle,draw=corPrim,minimum size=5pt,inner sep=0pt](b)at(1.05,.2){};
\node[circle,draw=corPrim,minimum size=5pt,inner sep=0pt](c)at(0,-.85){};
\node[anchor=north]at(-1.05,.04){L};
\node[anchor=north]at(1.05,.04){N};
\node[anchor=north]at(0,-1.04){PE};
\draw[thick](-3.4,.2)--(a);
\draw[thick](b)--(3.4,.2);
\draw[thick](0,-2.45)--(c);
\end{tikzpicture}
\column{.46\textwidth}
\begin{caixaProjeto}[title=O que o aluno deve enxergar]
A figura deixa separados os três bornes do mecanismo e evita a ambiguidade criada por um único retângulo representando toda a tomada.
\end{caixaProjeto}
\end{columns}
\begin{caixaObs}
Frente e traseira do mecanismo são vistas diferentes; a identificação impressa no componente deve aparecer no desenho didático.
\end{caixaObs}
\end{frame}

\begin{frame}{Várias tomadas — mostrar cada ponto como derivação}
\centering
\begin{tikzpicture}[scale=.79,every node/.style={font=\scriptsize}]
\node[draw=corPrim,fill=corDef!5,rounded corners,minimum width=1.9cm,minimum height=.85cm](q)at(-5.4,0){origem};
\foreach \x/\n in {-2.6/1,.6/2,3.8/3}{
  \node[draw=corNorma,fill=corNorma!5,rounded corners,minimum width=2.1cm,minimum height=1.15cm](t\n)at(\x,0){T\n\\L \quad N \quad PE};
}
\draw[very thick,corPrim!25](q)--(t1)--(t2)--(t3);
\draw[corPrim,thick]($(q.east)+(0,.28)$)--($(t1.west)+(0,.28)$)--($(t1.east)+(0,.28)$)--($(t2.west)+(0,.28)$)--($(t2.east)+(0,.28)$)--($(t3.west)+(0,.28)$);
\draw[corDef,thick](q.east)--(t1.west)--(t1.east)--(t2.west)--(t2.east)--(t3.west);
\draw[corNorma,thick]($(q.east)+(0,-.28)$)--($(t1.west)+(0,-.28)$)--($(t1.east)+(0,-.28)$)--($(t2.west)+(0,-.28)$)--($(t2.east)+(0,-.28)$)--($(t3.west)+(0,-.28)$);
\end{tikzpicture}
\begin{caixaObs}
O objetivo gráfico é mostrar três caminhos contínuos e três pontos de utilização independentes, em vez de caixas ligadas por uma seta genérica.
\end{caixaObs}
\end{frame}
